\begin{claim}{B.34}{Proved}{The complete genus-three subcatalogue containing a hyperelliptic involution}
There are exactly $79$ conjugacy classes of finite subgroups of the
extended mapping class group of a closed orientable genus-three surface
that contain a hyperelliptic involution.  Exactly $18$ preserve orientation
and $61$ contain orientation-reversing elements.  The latter $61$ classes
are listed below and are specified, together with every section and its
extension incidences, in
\texttt{type\_3\_0\_hyperelliptic\_mixed\_catalogue.csv} and
\texttt{type\_3\_0\_hyperelliptic\_sections.json}.
\end{claim}
\begin{proof}
Fix a hyperelliptic involution $z$ and its branched quotient
$p:S_3\longrightarrow S^2$, with branch set $P$ of size eight.
On the unbranched complements this is the connected double cover
classified by the homomorphism
\[
 \pi_1(S^2\setminus P)\longrightarrow C_2
 \quad\hbox{sending each of the eight meridians to }1.
                                                        \tag{B.34.1}
\]
Their sum is zero modulo two.  Every homeomorphism of the sphere
preserving $P$ preserves (B.34.1), and therefore has exactly two lifts.
The Birman--Hilden isotopy property gives the exact sequence
\[
 1\longrightarrow\langle z\rangle\longrightarrow
 C_{\mathcal T_{3,0}}(z)\longrightarrow\mathcal T_{0,8}
 \longrightarrow1.                                    \tag{B.34.2}
\]
Here the left group is the centralizer of this specified involution;
the map is induced by the quotient $p$.  The lifting and isotopy argument
is precisely the fully ramified hyperelliptic case of
\cite[\S9.4]{bFarbMargalit2012}.  Allowing orientation reversal adds the
negative lifts of the same sphere homeomorphisms, with the same kernel.

The involution $z$ is the unique hyperelliptic involution inside any
finite subgroup $K$ containing it.  To verify the assertion directly,
realize $K$ by isometries and use the resulting invariant complex
structure.  A hyperelliptic involution acts as $-1$ on $H^1(S_3;\mathbb R)$:
its invariant subspace is the pullback of $H^1(S^2;\mathbb R)=0$ under
its quotient, as follows by averaging cochains, or by the transfer
outside the branch points and filling their meridians.  If $z'$ were
another, the finite-order holomorphic map $zz'$ would act trivially on
$H^1$.  A nonidentity such map has isolated fixed points, each of
positive local Lefschetz index; a fixed-point-free map has Lefschetz
number zero.  But its Lefschetz number from homology is $2-6=-4$.
Thus $zz'=1$, proving uniqueness.
Consequently every element of $K$ centralizes $z$, and any conjugacy
between two such finite groups maps their hyperelliptic involutions
to one another.  Once $z$ is fixed, (B.34.2) therefore gives a bijection
\[
 \left\{\begin{array}{c}\text{finite subgroup conjugacy classes}\\
 \text{in }\mathcal T_{3,0}\text{ containing }z\end{array}\right\}
 \longleftrightarrow
 \left\{\begin{array}{c}\text{finite subgroup conjugacy classes}\\
 \text{in }\mathcal T_{0,8}\end{array}\right\}.
                                                        \tag{B.34.3}
\]
The inverse takes the full preimage, of order twice the sphere group.
The inclusion of every kernel element follows from the assumption
$z\in K$.  This also proves that no additional fusion occurs when the
chosen hyperelliptic involution is forgotten.

For completeness, B.26's decorated spherical-orbifold enumeration
applies verbatim to eight marked points: a selected cone orbit has
size $|H|/m$, a selected corner orbit has size $|H|/(2m)$, an orbit
on an open mirror edge has size $|H|/2$, and a regular orbit has size
$|H|$.  Their sum is eight.  The cyclic and dihedral rotation orders
are at most eight, since every non-polar point orbit has at least that
many points.  The tetrahedral and octahedral cases are included.
An icosahedral orbit has at least twelve points, so cannot occur.
The spherical classification and the homeomorphism classification
of these decorations in B.26 prove exhaustiveness and irredundancy.
Executing the integer weight equation gives $79$ decorated quotients,
$18$ spherical and $61$ with a negative element.

For each decorated quotient, the full lift is given by the exact
central extension presentation (B.28.3)--(B.28.5), now with eight
branch points.  Its presentation proof uses (B.34.1) in exactly the
same way: the displayed group surjects onto the full geometric lift
group of order $2|H|$, while its quotient by the central element $z$
is $H$, so it has at most $2|H|$ elements.  The two groups are
therefore isomorphic by that surjection.  For a projective-plane
quotient its generators include both $z$ and the crosscap element;
their relation is $u^2\prod x_i=z^a$.  In particular no cyclicity
assumption from the six-branch-point case is imposed.

The program \texttt{catalogue\_genus3\_hyperelliptic.py} was executed.
For every presentation it computes an exact regular permutation
representation of degree $2|H|$, verifies every defining relator,
checks that $z$ has order two, and checks injectivity on every local
cyclic or dihedral stabilizer.  The resulting orbifold Euler
characteristic is $-4/|K|$, so its oriented surface kernel has genus
three.  Exactly eight point sections have image $\langle z\rangle$,
providing an independent branch-number check.
The finite coset incidence graph of B.26 then computes every point
and oval fixed component and every partial section.  Each point
section has trivial centralizer in the surface group, each oval
section has infinite-cyclic centralizer, and the rotation multipliers,
extensions, maximal sections $\Phi$, group orbits and stabilizers
are recorded explicitly in the JSON file.  The calculation retains
all intersections of reflection ovals at dihedral stabilizers.
\end{proof}

In the following table, $b$ is both the G3H class suffix and the decorated sphere-eight row number
in \texttt{type\_3\_0\_hyperelliptic\_catalogue.csv}; $N$ is the group
order; $p$ counts all nontrivial partial splitting classes; and $m,o$
are respectively $|\Phi|$ and $|K\backslash\Phi|$.  The CSV gives the
cone, corner, mirror and ordinary branch decorations of every row,
and the JSON gives its finite group generators and all relators.
\begin{longtable}{r|r|r|r|r}
b&N&p&m&o\\\hline
\endhead
2&4&8&8&4\\
3&4&10&10&5\\
4&4&12&8&5\\
5&4&16&8&6\\
6&4&20&8&7\\
7&4&24&8&8\\
11&8&12&12&3\\
12&8&14&14&4\\
13&8&20&12&4\\
14&8&28&12&5\\
15&8&14&8&3\\
16&8&22&8&4\\
17&8&20&12&4\\
18&8&26&12&5\\
19&8&34&12&6\\
20&8&24&12&5\\
21&8&32&12&6\\
22&8&32&12&6\\
23&8&24&8&4\\
24&8&32&8&5\\
25&8&32&8&5\\
29&12&12&8&2\\
30&12&14&10&3\\
31&12&24&8&3\\
32&12&30&8&3\\
33&12&42&8&4\\
34&12&42&8&4\\
39&16&16&12&2\\
40&16&28&20&3\\
41&16&40&20&4\\
42&16&28&12&3\\
43&16&18&14&3\\
44&16&32&12&3\\
45&16&38&20&4\\
46&16&52&20&5\\
47&16&50&20&5\\
48&16&48&12&4\\
49&16&48&12&4\\
50&16&36&12&3\\
51&16&48&12&4\\
52&16&44&12&4\\
56&24&54&20&3\\
57&24&36&8&2\\
58&24&28&8&2\\
59&24&68&20&4\\
60&24&66&20&4\\
61&24&66&8&3\\
62&24&72&8&3\\
65&28&45&10&3\\
66&32&60&28&3\\
67&32&44&12&2\\
68&32&36&12&2\\
69&32&82&28&4\\
70&32&84&28&4\\
71&32&76&12&3\\
72&32&72&12&3\\
74&48&54&20&2\\
75&48&120&20&3\\
76&48&120&20&3\\
78&64&136&28&3\\
79&96&210&44&3\\
\end{longtable}
The mixed-class order distribution is
\[
\begin{array}{c|rrrrrrrrrr}
 |K|&4&8&12&16&24&28&32&48&64&96\\\hline
 \#&6&15&6&14&7&1&7&3&1&1.
\end{array}                                             \tag{B.34.4}
\]
The scope of (B.34.3) is the condition that the subgroup itself
contain a hyperelliptic involution.  A finite subgroup that only
centralizes an involution outside itself is represented by a
complement in a lift group; its conjugacy after forgetting the
external involution requires an additional calculation.  No such
complement is counted among these $79$ classes.
